% !TEX root = ../main.tex
\begin{ouverture}[fluides]{17}{Écoulement des fluides}{Partie 2 · Mouvement et interactions}
{Poussée d'Archimède\enspace·\enspace Flottaison\enspace·\enspace Débit volumique\enspace·\enspace Relation de Bernoulli\enspace·\enspace Effet Venturi}
\objectif{Expliquer qualitativement l'origine de la poussée d'Archimède et utiliser son expression vectorielle.}
\objectif{Prévoir si un corps flotte et exploiter la condition d'équilibre d'un corps flottant.}
\objectif{Exploiter la conservation du débit volumique pour déterminer la vitesse d'un fluide incompressible.}
\objectif{Exploiter la relation de Bernoulli, fournie, pour étudier qualitativement puis quantitativement un écoulement
permanent.}
\objectif{Modéliser la vidange d'un réservoir en Python.}
\prerequis{Pression, loi fondamentale de la statique des fluides (classe de première).}
\prerequis{Deuxième loi de Newton, équilibre d'un système (chapitre~14).}
\prerequis{Énergie mécanique et sa conservation (chapitre~15).}
\end{ouverture}
\essentiel
\essSection{1}{La poussée d'Archimède}
\begin{propriete}[Origine de la poussée d'Archimède]
Un corps plongé dans un fluide au repos subit, sur toute sa surface, des forces pressantes perpendiculaires à la paroi.
D'après la loi fondamentale de la statique des fluides, $p_\mathrm{B}-p_\mathrm{A}=\rho_\mathrm{f}\,g\,(z_\mathrm{A}-z_\mathrm{B})$ :
la pression augmente avec la profondeur. Les forces pressantes exercées sur la partie basse du corps l'emportent donc sur
celles exercées sur sa partie haute, tandis que les forces latérales se compensent : leur résultante est verticale et
dirigée vers le haut.
\end{propriete}
\schema{17-archimede}
\begin{loi}[Poussée d'Archimède]
Tout corps immergé dans un fluide au repos, de masse volumique $\rho_\mathrm{f}$, subit de la part de ce fluide une force
verticale, dirigée vers le haut, égale à l'opposé du poids du fluide déplacé :
\[\vec{\Pi}=-\rho_\mathrm{f}\,V_\mathrm{imm}\,\vec{g}\qquad\text{soit}\qquad \Pi=\rho_\mathrm{f}\,V_\mathrm{imm}\,g,\]
avec $\Pi$ en newtons, $\rho_\mathrm{f}$ en \unit{kg.m^{-3}} et $V_\mathrm{imm}$, volume de la partie immergée, en \unit{m^3}.
\end{loi}
\begin{savoirfaire}[Retrouver l'expression pour un pavé immergé]
Pour un pavé de hauteur $h$ et de base $S$, entièrement immergé, la différence de pression entre les faces inférieure
et supérieure vaut $p_\mathrm{bas}-p_\mathrm{haut}=\rho_\mathrm{f}\,g\,h$. La résultante des forces pressantes est donc
verticale, vers le haut, de valeur
\[F_\mathrm{bas}-F_\mathrm{haut}=(p_\mathrm{bas}-p_\mathrm{haut})\,S=\rho_\mathrm{f}\,g\,h\,S=\rho_\mathrm{f}\,V\,g.\]
Elle ne dépend pas de la profondeur du pavé, mais seulement de son volume.
\end{savoirfaire}
\begin{propriete}[Couler ou flotter]
Un corps homogène de masse volumique $\rho$, lâché entièrement immergé, est soumis à son poids $P=\rho\,V\,g$ et à la
poussée $\Pi=\rho_\mathrm{f}\,V\,g$. Si $\rho>\rho_\mathrm{f}$, il coule ; si $\rho<\rho_\mathrm{f}$, il remonte et flotte.
À l'équilibre, la poussée compense le poids, $\rho_\mathrm{f}\,V_\mathrm{imm}\,g=\rho\,V\,g$, d'où
\[\frac{V_\mathrm{imm}}{V}=\frac{\rho}{\rho_\mathrm{f}}.\]
\end{propriete}
\begin{exemple}
La glace a pour masse volumique \qty{917}{kg.m^{-3}}. Dans l'eau de mer ($\rho_\mathrm{f}=\qty{1025}{kg.m^{-3}}$), la
fraction immergée d'un iceberg vaut $\dfrac{917}{1025}\approx\qty{89}{\percent}$ : seuls 11~\% de son volume émergent.
Un glaçon dans un verre d'eau douce est immergé à \qty{91,7}{\percent}.
\end{exemple}
\essSection{2}{Écoulement d'un fluide en régime permanent}
\begin{definition}[Régime permanent et débit volumique]
L'écoulement d'un fluide est \textbf{permanent} (ou stationnaire) si, en chaque point, la vitesse et la pression du
fluide ne dépendent pas du temps. Le \textbf{débit volumique} $D_V$ est le volume de fluide qui traverse une section par
unité de temps ; si le fluide traverse la section $S$ à la vitesse moyenne $v$, perpendiculaire à $S$,
\[D_V=\frac{V}{\Delta t}=S\,v\qquad(D_V\text{ en }\unit{m^3.s^{-1}},\ S\text{ en }\unit{m^2},\ v\text{ en }\unit{m.s^{-1}}).\]
\end{definition}
\begin{propriete}[Conservation du débit volumique]
Un liquide se comporte comme un fluide \textbf{incompressible} : sa masse volumique reste constante. En régime permanent,
le débit volumique se conserve tout au long d'une conduite, même si sa section varie :
\[S_\mathrm{A}\,v_\mathrm{A}=S_\mathrm{B}\,v_\mathrm{B}.\]
Là où la section diminue, la vitesse augmente.
\end{propriete}
\begin{exemple}
Le cœur envoie \qty{5,0}{L} de sang par minute, soit $D_V=\qty{8,3e-5}{m^3.s^{-1}}$, dans l'aorte, de diamètre
\qty{2,5}{cm} : le sang y circule à $v=D_V/S=\qty{0,17}{m.s^{-1}}$. Ce débit se répartit ensuite dans des milliards de
capillaires dont la section totale atteint environ \qty{0,25}{m^2} : la vitesse n'y vaut plus que \qty{0,33}{mm.s^{-1}},
ce qui laisse le temps aux échanges avec les cellules.
\end{exemple}
\essSection{3}{La relation de Bernoulli}
\begin{loi}[Relation de Bernoulli]
Pour un fluide incompressible, en écoulement permanent, dont on néglige les frottements, la grandeur
\[p+\frac12\,\rho\,v^2+\rho\,g\,z\]
garde la même valeur en tout point d'une même ligne de courant. Chaque terme est homogène à une pression, c'est-à-dire
à une énergie par unité de volume : la relation traduit la conservation de l'énergie du fluide. Elle est fournie dans
les exercices.
\end{loi}
\begin{propriete}[Effet Venturi]
Dans une conduite horizontale, là où la section diminue, la vitesse augmente (conservation du débit) et la pression
\textbf{diminue} (relation de Bernoulli à $z$ constant) :
\[p_\mathrm{A}-p_\mathrm{B}=\frac12\,\rho\,\big(v_\mathrm{B}^2-v_\mathrm{A}^2\big).\]
Cet effet est exploité dans les trompes à eau, les pulvérisateurs, les débitmètres et les carburateurs.
\end{propriete}
\schema{17-venturi}
\begin{savoirfaire}[Établir la formule de Torricelli]
Un grand réservoir ouvert se vide par un petit orifice B situé à la profondeur $h$ sous la surface libre A. En A et en
B, le fluide est à la pression atmosphérique ; la surface libre descend très lentement, $v_\mathrm{A}\approx0$. La relation
de Bernoulli entre A et B s'écrit $p_\mathrm{atm}+\rho\,g\,z_\mathrm{A}=p_\mathrm{atm}+\frac12\,\rho\,v_\mathrm{B}^2+\rho\,g\,z_\mathrm{B}$,
d'où
\[v_\mathrm{B}=\sqrt{2\,g\,h}.\]
Pour $h=\qty{1,25}{m}$, l'eau sort à \qty{4,95}{m.s^{-1}}, comme si elle était tombée en chute libre de la hauteur $h$.
\end{savoirfaire}
\schema{17-torricelli}
\begin{codecote}[0.62]
\begin{pythoncol}
from math import pi, sqrt
g = 9.81
S = pi * 0.050**2       # récipient (m²)
s = pi * 0.002**2       # orifice (m²)
h, t, dt = 0.20, 0.0, 0.01
while h > 0:
    v = sqrt(2 * g * h)       # Torricelli
    h = h - s / S * v * dt    # débit conservé
    t = t + dt
T = S / s * sqrt(2 * 0.20 / g)   # exacte
print(round(t, 1), round(T, 1))
\end{pythoncol}
\begin{consolecol}
126.2 126.2
\end{consolecol}
\cote
\fontsize{9.5}{12.5}\selectfont
\capacite{Capacité numérique}\par\smallskip
Un récipient cylindrique de \qty{10}{cm} de diamètre, rempli sur \qty{20}{cm}, se vide par un orifice de \qty{4}{mm} de
diamètre. Pendant $\mathrm{d}t$, le volume $S\,|\mathrm{d}h|$ sort à la vitesse de Torricelli : $S\,|\mathrm{d}h|=s\,v\,\mathrm{d}t$.
La méthode d'Euler donne une durée de vidange de \qty{126,2}{s}, en accord avec l'expression exacte
$\dfrac{S}{s}\sqrt{\dfrac{2H}{g}}$.
\end{codecote}
\begin{attention}
La poussée d'Archimède dépend du volume immergé et de la masse volumique du \textbf{fluide}, pas de la masse volumique du
corps, ni de sa profondeur une fois qu'il est entièrement immergé. Dans la relation de Bernoulli, $p$ s'exprime en
pascals, $\rho$ en \unit{kg.m^{-3}}, $v$ en \unit{m.s^{-1}} et $z$ en mètres, l'axe $(Oz)$ étant vertical vers le haut.
\end{attention}
\begin{experience}[Tester l'expression de la poussée d'Archimède]
On suspend un cylindre à un dynamomètre, puis on l'immerge progressivement dans un liquide. La différence entre le poids
mesuré dans l'air et l'indication du dynamomètre donne la poussée. Tracée en fonction du volume immergé, elle forme une
droite passant par l'origine, de coefficient directeur $\rho_\mathrm{f}\,g$ ; elle ne varie plus une fois le cylindre
entièrement immergé. Dans de l'eau salée, on obtient \qty{10,1}{N} par litre immergé, soit
$\rho_\mathrm{f}=\qty{1,03e3}{kg.m^{-3}}$.
\end{experience}
\begin{rediger}
Avant d'utiliser la relation de Bernoulli, rappeler ses conditions : fluide incompressible, écoulement permanent,
frottements négligés. Choisir deux points A et B d'une même ligne de courant et préciser ce que l'on sait en chacun :
« à l'air libre, $p=p_\mathrm{atm}$ » ; « surface libre d'un grand réservoir, $v\approx0$ ».
\end{rediger}
\exercices

\rubrique{Découverte}
\begin{blocLivre}{ex-17-1}
\begin{exercice}[1]{D'où vient la poussée d'Archimède ?}
Un cube d'arête $a=\qty{10,0}{cm}$ est maintenu immergé dans l'eau, ses faces horizontales, sa face supérieure à
\qty{20,0}{cm} sous la surface libre.
\begin{donnees}
$\rho_\mathrm{eau}=\qty{1000}{kg.m^{-3}}$ ; $g=\qty{9,81}{N.kg^{-1}}$ ; $p_\mathrm{atm}=\qty{1,013e5}{Pa}$.
\end{donnees}
\begin{questions}
\item Calculer la pression de l'eau sur la face supérieure, puis sur la face inférieure du cube.
\item En déduire les valeurs des forces pressantes exercées sur ces deux faces.
\item Pourquoi les forces pressantes exercées sur les faces latérales se compensent-elles ? Calculer la valeur de la
résultante des forces pressantes et la comparer à $\rho_\mathrm{eau}\,V\,g$.
\item Le cube est descendu \qty{1}{m} plus bas. La résultante change-t-elle ? Justifier.
\end{questions}
\end{exercice}
\end{blocLivre}
\begin{blocLivre}{ex-17-2}
\begin{exercice}[1]{La couronne du roi Hiéron}
Selon la légende, le roi Hiéron de Syracuse demanda à Archimède de vérifier, sans l'abîmer, que sa couronne était en or
pur. Suspendue à un dynamomètre, une couronne de masse \qty{1,000}{kg} donne l'indication \qty{9,810}{N} dans l'air et
\qty{9,180}{N} lorsqu'elle est entièrement immergée dans l'eau.
\begin{donnees}
$\rho_\mathrm{eau}=\qty{1000}{kg.m^{-3}}$ ; $g=\qty{9,81}{N.kg^{-1}}$ ; masse volumique de l'or : \qty{19300}{kg.m^{-3}}.
\end{donnees}
\begin{questions}
\item Faire le bilan des forces exercées sur la couronne immergée et calculer la poussée d'Archimède.
\item En déduire le volume de la couronne.
\item Calculer sa masse volumique. La couronne est-elle en or pur ?
\item Quelle indication le dynamomètre aurait-il donnée pour une couronne d'or pur de même masse ?
\end{questions}
\end{exercice}
\end{blocLivre}
\begin{blocLivre}{ex-17-3}
\begin{exercice}[1]{Le tuyau d'arrosage}
Un tuyau d'arrosage remplit un seau de \qty{10,0}{L} en \qty{50}{s}. Son diamètre intérieur vaut \qty{15}{mm} ; on
visse à son extrémité un embout de \qty{5,0}{mm} de diamètre.
\begin{questions}
\item Calculer le débit volumique, en \unit{m^3.s^{-1}} puis en litres par minute.
\item Calculer la vitesse de l'eau dans le tuyau.
\item En utilisant la conservation du débit, calculer la vitesse de l'eau à la sortie de l'embout. Par combien
a-t-elle été multipliée ?
\item À quelle hauteur monterait un jet vertical sortant de l'embout, frottements négligés ? Pourquoi, sans embout,
pince-t-on l'extrémité du tuyau pour arroser plus loin ?
\end{questions}
\end{exercice}
\end{blocLivre}
\begin{blocLivre}{ex-17-4}
\begin{exercice}[1]{L'effet Venturi}
De l'eau s'écoule dans une conduite horizontale qui comporte un étranglement, comme sur le schéma de l'essentiel. La
section de l'étranglement est quatre fois plus petite que celle de la conduite, où l'eau circule à
\qty{0,50}{m.s^{-1}}.
\begin{donnees}
$\rho_\mathrm{eau}=\qty{1000}{kg.m^{-3}}$ ; $g=\qty{9,81}{m.s^{-2}}$ ; relation de Bernoulli :
$p+\frac12\,\rho\,v^2+\rho\,g\,z=\text{constante}$.
\end{donnees}
\begin{questions}
\item Dans quelle partie la vitesse est-elle la plus grande ? Calculer la vitesse dans l'étranglement.
\item Dans quelle partie la pression est-elle la plus faible ? Calculer la différence de pression entre la conduite et
l'étranglement.
\item Dans les tubes verticaux, la différence de hauteur d'eau vérifie $p_\mathrm{A}-p_\mathrm{B}=\rho\,g\,\Delta h$.
Calculer $\Delta h$.
\item Lorsqu'on souffle entre deux feuilles de papier tenues verticalement, elles se rapprochent. Expliquer.
\end{questions}
\schema{17-venturi}
\end{exercice}
\end{blocLivre}
\begin{blocLivre}{ex-17-5}
\begin{exercice}[1]{Le tonneau percé}
Un tonneau rempli d'eau, ouvert à l'air libre, est percé d'un orifice circulaire de \qty{1,0}{cm} de diamètre, situé
\qty{0,80}{m} sous la surface libre et \qty{0,50}{m} au-dessus du sol. Le jet sort horizontalement.
\begin{donnees}
$g=\qty{9,81}{m.s^{-2}}$ ; formule de Torricelli : $v=\sqrt{2\,g\,h}$.
\end{donnees}
\begin{questions}
\item Calculer la vitesse de l'eau à la sortie de l'orifice.
\item En déduire le débit volumique, en litres par seconde puis en litres par minute.
\item En utilisant les résultats du chapitre~15, calculer la distance entre le pied du tonneau et le point où le jet
touche le sol.
\item Comment évoluent la vitesse du jet et cette distance au cours de la vidange ? Justifier.
\end{questions}
\end{exercice}
\end{blocLivre}
\rubrique{Entraînement}
\begin{blocLivre}{ex-17-6}
\begin{exercice}[2]{Le ballon-sonde}
Météo-France lâche chaque jour des ballons-sondes gonflés à l'hélium, qui emportent des instruments de mesure. Au sol,
le ballon a un volume de \qty{3,0}{m^3} ; l'enveloppe et les instruments ont une masse totale de \qty{1,50}{kg}.
\begin{donnees}
Au sol : $\rho_\mathrm{air}=\qty{1,20}{kg.m^{-3}}$ ; $\rho_\mathrm{He}=\qty{0,17}{kg.m^{-3}}$ ; $g=\qty{9,81}{N.kg^{-1}}$.
\end{donnees}
\begin{questions}
\item Faire le bilan des forces exercées sur le système \{ballon, hélium, instruments\} juste après le lâcher, lorsque
sa vitesse est encore nulle.
\item Calculer la poussée d'Archimède et le poids du système.
\item Appliquer la deuxième loi de Newton pour calculer l'accélération du ballon au moment du lâcher.
\item En altitude, la pression de l'air diminue. Comment évolue le volume du ballon ? Pourquoi finit-il par éclater,
vers \qty{30}{km} d'altitude ?
\end{questions}
\end{exercice}
\end{blocLivre}
\begin{blocLivre}{ex-17-7}
\begin{exercice}[2]{Une artère rétrécie}
Une plaque d'athérome réduit de moitié le diamètre d'une artère, sur quelques millimètres : c'est une sténose. En
amont, l'artère a un diamètre de \qty{6,0}{mm} et le sang y circule à \qty{0,40}{m.s^{-1}}. On assimile le sang à un fluide
incompressible, en écoulement permanent, et l'on néglige les frottements.
\begin{donnees}
$\rho_\mathrm{sang}=\qty{1060}{kg.m^{-3}}$ ; $\qty{1}{mmHg}=\qty{133,3}{Pa}$ ; relation de Bernoulli :
$p+\frac12\,\rho\,v^2+\rho\,g\,z=\text{constante}$.
\end{donnees}
\begin{questions}
\item Calculer le débit sanguin dans l'artère, en \unit{m^3.s^{-1}} puis en litres par minute.
\item Calculer la vitesse du sang dans la sténose.
\item L'artère étant horizontale, calculer la chute de pression dans la sténose, en pascals puis en millimètres de
mercure.
\item Pourquoi une sténose importante risque-t-elle de faire s'affaisser la paroi de l'artère ? Le médecin entend alors
un « souffle » au stéthoscope : quelle est son origine ?
\end{questions}
\end{exercice}
\end{blocLivre}
\begin{blocLivre}{ex-17-8}
\begin{exercice}[2]{Le débitmètre de Venturi}
Pour mesurer le débit d'une canalisation horizontale, on y insère un tube de Venturi, comme sur le schéma de
l'essentiel : la section passe de $S_\mathrm{A}=\qty{20}{cm^2}$ à $S_\mathrm{B}=\qty{5,0}{cm^2}$. La différence des hauteurs
d'eau dans les tubes verticaux vaut $\Delta h=\qty{12,0}{cm}$.
\begin{donnees}
$\rho_\mathrm{eau}=\qty{1000}{kg.m^{-3}}$ ; $g=\qty{9,81}{m.s^{-2}}$ ; relation de Bernoulli :
$p+\frac12\,\rho\,v^2+\rho\,g\,z=\text{constante}$.
\end{donnees}
\begin{questions}
\item Exprimer $v_\mathrm{B}$ en fonction de $v_\mathrm{A}$.
\item L'eau étant immobile dans les tubes verticaux, justifier que $p_\mathrm{A}-p_\mathrm{B}=\rho\,g\,\Delta h$. En
déduire que $v_\mathrm{A}=\sqrt{\dfrac{2\,g\,\Delta h}{(S_\mathrm{A}/S_\mathrm{B})^2-1}}$.
\item Calculer $v_\mathrm{A}$, puis le débit, en litres par seconde et en \unit{m^3.h^{-1}}.
\item Montrer que le débit est proportionnel à $\sqrt{\Delta h}$. Quelle différence de hauteur lirait-on pour un débit
double ?
\end{questions}
\schema{17-venturi}
\end{exercice}
\end{blocLivre}
\begin{blocLivre}{ex-17-9}
\begin{exercice}[2]{Simuler la vidange d'un réservoir}
On reprend le programme de l'essentiel, qui simule la vidange d'un récipient cylindrique par la méthode d'Euler.
\begin{questions}
\item \capacite{Capacité numérique}Quelle relation traduit la ligne \code{v = sqrt(2 * g * h)} ? Justifier la ligne
\code{h = h - s / S * v * dt} par la conservation du débit.
\item Pourquoi la boucle s'arrête-t-elle ? Que représente la valeur de \code{t} affichée ? Comparer les deux nombres
affichés.
\item Modifier le programme pour enregistrer $t$ et $h$ dans des listes, puis tracer $h$ en fonction de $t$. Décrire
l'allure de la courbe obtenue.
\item Comment varie la durée de vidange si le diamètre de l'orifice double ? Vérifier avec le programme.
\item Une mesure donne \qty{152}{s}. Proposer une explication à l'écart observé.
\end{questions}
\end{exercice}
\end{blocLivre}
\begin{blocLivre}{ex-17-10}
\begin{exercice}[2]{Le densimètre}
Un densimètre est un tube cylindrique lesté, de masse $m=\qty{24,0}{g}$, de section $S=\qty{2,00}{cm^2}$ et de longueur
$L=\qty{15,0}{cm}$. Il flotte verticalement ; on lit la profondeur à laquelle il s'enfonce.
\schema{17-ex10}
\begin{donnees}
$\rho_\mathrm{eau}=\qty{1000}{kg.m^{-3}}$ ; $g=\qty{9,81}{N.kg^{-1}}$.
\end{donnees}
\begin{questions}
\item Écrire la condition d'équilibre du densimètre et montrer que la profondeur immergée vaut $h=\dfrac{m}{\rho\,S}$,
$\rho$ étant la masse volumique du liquide. Vérifier la valeur indiquée pour l'eau.
\item Calculer la masse volumique du liquide inconnu.
\item Quelle est la plus petite masse volumique mesurable avec ce densimètre ?
\item Calculer la position de la graduation \qty{1100}{kg.m^{-3}}. Les graduations sont-elles régulièrement espacées ?
\end{questions}
\end{exercice}
\end{blocLivre}
\begin{blocLivre}{ex-17-11}
\begin{exercice}[2]{Le château d'eau}
La surface libre de l'eau d'un château d'eau se trouve \qty{35}{m} au-dessus du robinet d'une cuisine, au rez-de-chaussée
d'un immeuble. Un autre robinet est situé au troisième étage, \qty{9,0}{m} plus haut.
\begin{donnees}
$\rho_\mathrm{eau}=\qty{1000}{kg.m^{-3}}$ ; $g=\qty{9,81}{m.s^{-2}}$ ; $p_\mathrm{atm}=\qty{1,013e5}{Pa}$ ; relation de Bernoulli :
$p+\frac12\,\rho\,v^2+\rho\,g\,z=\text{constante}$.
\end{donnees}
\begin{questions}
\item Les robinets étant fermés, l'eau est immobile. Calculer la pression de l'eau au robinet du rez-de-chaussée.
\item Calculer la pression au robinet du troisième étage. Pourquoi les châteaux d'eau sont-ils construits sur des points
hauts ?
\item On ouvre le robinet du rez-de-chaussée. En négligeant les frottements, calculer la vitesse de l'eau à la sortie.
\item Ce robinet, de section \qty{1,0}{cm^2}, débite en réalité \qty{12}{L} par minute. Calculer la vitesse réelle et
expliquer l'écart.
\end{questions}
\end{exercice}
\end{blocLivre}
\rubrique{Objectif bac}
\begin{blocLivre}{ex-17-12}
\begin{exercice}[3]{Le tube de Pitot}
\begin{docu}{1}{Mesurer la vitesse d'un avion}
Les avions mesurent leur vitesse par rapport à l'air grâce à des tubes de Pitot, fixés sous le nez de l'appareil. Un
capteur mesure la différence $\Delta p$ entre la pression au point d'arrêt B, où l'air est immobile par rapport à la
sonde, et la pression au niveau des orifices latéraux A, où l'air s'écoule à la vitesse $v$.
\end{docu}
\schema{17-ex12}
\begin{docu}{2}{Deux phases du vol}
En croisière, à \qty{11000}{m} d'altitude, $\rho_\mathrm{air}=\qty{0,41}{kg.m^{-3}}$ et le capteur indique
$\Delta p=\qty{12,8}{kPa}$. Au niveau du sol, $\rho_\mathrm{air}=\qty{1,22}{kg.m^{-3}}$ ; l'avion décolle à \qty{280}{km.h^{-1}}.
\end{docu}
\begin{donnees}
Relation de Bernoulli : $p+\frac12\,\rho\,v^2+\rho\,g\,z=\text{constante}$ ; on néglige la différence d'altitude entre A et B
et la compressibilité de l'air.
\end{donnees}
\begin{questions}
\item En appliquant la relation de Bernoulli entre un point de l'écoulement à la vitesse $v$ et le point d'arrêt, montrer
que $\Delta p=\frac12\,\rho_\mathrm{air}\,v^2$.
\item Calculer la vitesse de l'avion en croisière, en \unit{m.s^{-1}} puis en \unit{km.h^{-1}}.
\item Quelle différence de pression le capteur mesure-t-il au moment du décollage ?
\item \initiative En croisière, quelle vitesse afficherait un instrument qui utiliserait toujours la masse volumique de
l'air au sol ? Sachant que la portance des ailes est proportionnelle à $\frac12\,\rho_\mathrm{air}\,v^2$, expliquer pourquoi
cette « vitesse indiquée » reste très utile au pilote.
\end{questions}
\end{exercice}
\end{blocLivre}
\begin{blocLivre}{ex-17-13}
\begin{exercice}[3]{La lance à incendie}
\begin{docu}{1}{Le matériel}
Une lance à incendie est alimentée par un tuyau de \qty{70}{mm} de diamètre intérieur ; son embout de sortie a un
diamètre de \qty{22}{mm}. Le débit vaut \qty{500}{L} par minute.
\end{docu}
\begin{docu}{2}{Sur le terrain}
Les pompiers savent qu'un jet dirigé verticalement monte en pratique jusqu'à environ \qty{18}{m}, nettement moins
haut que ce que prévoit le modèle sans frottement.
\end{docu}
\begin{donnees}
$\rho_\mathrm{eau}=\qty{1000}{kg.m^{-3}}$ ; $g=\qty{9,81}{m.s^{-2}}$ ; relation de Bernoulli :
$p+\frac12\,\rho\,v^2+\rho\,g\,z=\text{constante}$.
\end{donnees}
\begin{questions}
\item Calculer la vitesse de l'eau dans le tuyau et à la sortie de la lance.
\item La lance est horizontale. Calculer la surpression $p-p_\mathrm{atm}$ nécessaire à l'entrée de la lance.
\item \initiative Calculer la hauteur maximale que pourrait atteindre un jet vertical en l'absence de frottements. Un
pompier peut-il, en théorie, atteindre le toit d'un immeuble de \qty{20}{m} ?
\item Proposer des explications à l'écart avec le document~2.
\end{questions}
\end{exercice}
\end{blocLivre}
\begin{blocLivre}{ex-17-14}
\begin{exercice}[3]{La montgolfière}
\begin{docu}{1}{Une enveloppe d'air chaud}
Une montgolfière a une enveloppe de \qty{2500}{m^3}, ouverte à sa base. La masse de l'enveloppe, de la nacelle, du
brûleur et des bouteilles de gaz vaut \qty{400}{kg}. Le brûleur maintient l'air intérieur à \qty{100}{\degreeCelsius}.
\end{docu}
\begin{docu}{2}{Masse volumique de l'air}
À la pression atmosphérique, la masse volumique de l'air vaut $\rho=\dfrac{353}{T}$, en \unit{kg.m^{-3}}, la température $T$
étant exprimée en kelvins ($T=\theta+273$).
\end{docu}
\begin{donnees}
$g=\qty{9,81}{N.kg^{-1}}$ ; masse moyenne d'un passager, pilote compris : \qty{80}{kg} ; on néglige le volume de la
nacelle et des passagers.
\end{donnees}
\begin{questions}
\item Faire le bilan des forces exercées sur le système \{montgolfière, air chaud, passagers\} et écrire la condition de
décollage.
\item Calculer la masse volumique de l'air extérieur à \qty{15}{\degreeCelsius} et celle de l'air chaud. Calculer la poussée
d'Archimède et le poids de l'air chaud.
\item \initiative Combien de personnes la montgolfière peut-elle emporter par une matinée à \qty{15}{\degreeCelsius} ?
\item Reprendre le calcul pour un après-midi d'été à \qty{30}{\degreeCelsius}. Pourquoi les vols ont-ils lieu tôt le matin
ou en fin de journée ?
\end{questions}
\end{exercice}
\end{blocLivre}
\begin{blocLivre}{ex-17-15}
\begin{exercice}[3]{La centrale de haute chute}
\begin{docu}{1}{Une conduite forcée}
Dans une centrale hydroélectrique de montagne, l'eau d'un lac de retenue descend dans une conduite forcée de
\qty{2,0}{m} de diamètre jusqu'à la turbine, située \qty{600}{m} sous la surface du lac. Le débit vaut
\qty{20}{m^3.s^{-1}}.
\end{docu}
\begin{docu}{2}{Puissance disponible}
La puissance hydraulique que l'eau peut céder à la turbine vaut $P=\rho\,g\,h\,D_V$, $h$ étant la hauteur de chute. La
turbine et l'alternateur ont un rendement global de 90~\%. Un foyer consomme en moyenne \qty{4,5}{MWh} d'électricité
par an.
\end{docu}
\begin{donnees}
$\rho_\mathrm{eau}=\qty{1000}{kg.m^{-3}}$ ; $g=\qty{9,81}{m.s^{-2}}$ ; $p_\mathrm{atm}=\qty{1,013e5}{Pa}$ ; relation de Bernoulli :
$p+\frac12\,\rho\,v^2+\rho\,g\,z=\text{constante}$.
\end{donnees}
\begin{questions}
\item Calculer la vitesse de l'eau dans la conduite forcée.
\item En négligeant les frottements, appliquer la relation de Bernoulli entre la surface du lac et l'entrée de la
turbine pour calculer la pression de l'eau à l'entrée de la turbine. Commenter l'importance du terme
$\frac12\,\rho\,v^2$.
\item Calculer la puissance hydraulique, puis la puissance électrique de la centrale.
\item \initiative La centrale fonctionne \num{3000} heures par an. Combien de foyers peut-elle alimenter ?
\end{questions}
\end{exercice}
\end{blocLivre}
